\documentclass{hogent-article}

\begin{document}
    Dit interview werd afgenomen met Mevr. Katrijn Van den Broeck (klinisch psycholoog en co-promotor) in het kader van de requirementsanalyse. Het interview volgde een semi-gestructureerd formaat waarbij onderstaande vragen als leidraad dienden voor het verzamelen van klinische inzichten en functionele behoeften voor de applicatie.
    
    \vspace{1em}
    
    \begin{description}
        \item[De doelgroep \& probleemstelling] \hfill
        \begin{enumerate}
            \item Welke specifieke subgroepen ziet u binnen de twintigers (studenten, pas afgestudeerden, personen die al langer werken, \dots) en verschillen hun behoeften op vlak van ondersteuning bij burn-out?
            \item Heeft u al gebruik gemaakt van tools of apps die cliënten ondersteuning biedt voor het behandelen of voorkomen van burn-out? Met andere woorden: hoe kunnen ze zelf aan de slag buiten het raadplegen van psychologische sessies?
            \item Tijdens de literatuurstudie werd besproken hoe de huidige mentaal welzijn apps kampen met een zekere retentieproblematiek. Wat is volgens u de grootste reden waarom jongeren stoppen met het gebruiken van mentale welzijn-apps?
        \end{enumerate}
        
        \vspace{1em}
        
        \item[Werkwijze en methodiek] \hfill
        \begin{enumerate}
            \setcounter{enumi}{3}
            \item Welke specifieke behandelingen/methodieken past u zelf toe ter behandeling van burn-out?
            \item Hoe past u deze toe (bv. templates, kaartjes, dagboeken, fysieke notities, \dots)? Met andere woorden: hoe moet de praktische vertaling van deze methodiek eruitzien voor een app-ontwerp?
            \item Welke elementen van REBT, mindfulness of psycho-educatie lenen zich volgens u het beste voor een zelfhulp-app?
            \item REBT draait om het uitdagen van gedachten. Hoe ziet u dit idealiter in een app (bv. chat-vorm, invulformulier, dagelijkse reflectie, \dots)?
            \item Communities kunnen bijdragen aan sociale ondersteuning. Wat is uw visie op het gebruik van een community in combinatie met het uitoefenen van REBT-technieken?
        \end{enumerate}
        
        \vspace{1em}
        
        \item[Design \& stressfactoren (UI/UX)] \hfill
        \begin{enumerate}
            \setcounter{enumi}{8}
            \item Welke visuele elementen of functies (zoals notificaties) bij digitale toepassingen veroorzaken volgens u bij cliënten juist meer stress?
            \item Wat is de grootste 'no-go' in een app voor iemand die reeds burn-out symptomen vertoont (denk aan: hoeveelheid tekst, kleurgebruik, keuzestress)?
            \item Zijn er bepaalde kleuren of visuele structuren die u in de praktijk ziet die cliënten met burn-out symptomen direct afstoten of juist kalmeren?
            \item Hoe kunnen we groei of progressie binnen de app visueel maken zonder dat de gebruiker het gevoel krijgt dat hij 'moet presteren' of een 'streak' moet hooghouden?
            \item Waar ligt volgens u voor deze doelgroep de grens tussen een behulpzame herinnering en een stress-opwekkende 'technostressor'?
        \end{enumerate}
        
        \vspace{1em}
        
        \item[Implementatie \& functionaliteiten] \hfill
        \begin{enumerate}
            \setcounter{enumi}{13}
            \item Gaat uw voorkeur uit naar een applicatie gericht op preventie of op behandeling, en hoe zou dit in implementatie vorm moeten krijgen?
            \item Welke functionaliteiten zijn volgens u absoluut noodzakelijk (\textit{must-have}) om te implementeren in een applicatie om burn-out te bestrijden of te voorkomen?
        \end{enumerate}
    \end{description}
\end{document}